\documentclass[10pt,a4paper]{amsart}
\usepackage[margin=19mm]{geometry}
\usepackage{amsmath,amssymb,mathtools}
\usepackage[scr=rsfs,cal=euler]{mathalfa}
\usepackage{xfrac}
\usepackage[unicode,hidelinks,bookmarks=false]{hyperref}
\newcommand{\ie}{i.e.\ }
\newcommand{\cf}{cf.\ }
\newcommand{\resp}{respectively\ }
\newcommand{\Subsection}{\subsection}
\providecommand{\nobreakdash}{\nobreak\hbox{-}}
\setlength{\emergencystretch}{2em}
\multlinegap0pt
\usepackage{amssymb}
\newtheorem{theorem}{Theorem}
\title{On Numerical Semigroups with Fixed Quotient}
\author{Ignacio Ojeda, Jos\'e Carlos Rosales}
\date{Source: arXiv:2605.14775v1 (CC BY 4.0)}
\begin{document}
\maketitle

\noindent\emph{Source and licence.} On Numerical Semigroups with Fixed Quotient. Version arXiv:2605.14775v1 (CC BY 4.0), distributed by arXiv under the Creative Commons Attribution 4.0 International licence, http://creativecommons.org/licenses/by/4.0/, as recorded in the arXiv licence metadata for this exact version and shown on the abstract page https://arxiv.org/abs/2605.14775v1. Full licence notice: \texttt{licenses/arxiv-2605.14775-CC-BY-4.0.txt}.\par\smallskip

\noindent\textit{Editorial scope.} Counted unit: the article's theorem th:PF-rankone (source0.tex lines 979--985). The article's complete original proof of it is delivered with it (source0.tex lines 987--1035, one \texttt{proof} environment of 49 lines). No mathematics is written, repaired or re-derived: the statement and the proof are the article's own lines, in the article's own notation, and no retained line is substituted or re-flowed.  No further source-local context is needed: every cross-reference of every retained line resolves inside the two delivered spans.  Nothing was omitted from any retained span, so the package carries no omission record.  No retained line contains a citation, so the wrapper carries no bibliography and no bibliography entry.  No retained line contains a figure environment, an image-inclusion command or drawing code, so no image is delivered and the package declares no asset dependency.  Editorial formatting only, in addition: an AMS \texttt{amsart} preamble that restates the article's own declarations and macros used by the retained lines, namely the environments \texttt{theorem} on separate counters, together with the standard packages those lines need.  The retained environments print in this document as Theorem 1 in order of appearance, whereas the article prints the same items with the numbering of its own full document.  The complete native source of the article as distributed in the arXiv e-print for this exact version is bundled unchanged as \texttt{sources/200611/source0.tex}.
\begin{theorem}\label{th:PF-rankone}
Let $x\in \Delta$ with $\gcd(x,d)=1$, and set $S=\langle x\rangle+d\Delta$. Then $\operatorname{t}(S)=\operatorname{t}(\Delta)$ and
\[
\operatorname{PF}(S)=\{df+(d-1)x : f\in \operatorname{PF}(\Delta)\}.
\]
In particular, $S$ is symmetric if and only if $\Delta$ is symmetric.
\end{theorem}

\begin{proof}
Set $S=\langle x\rangle+d\Delta=\{ax+d\delta:\ a\in\mathbb{N},\ \delta\in\Delta\}$.
For every \(k\in\mathbb{Z}\), we first record the equivalence
\[
(d-1)x+dk\in S \iff k\in\Delta.
\]
If $k\in\Delta$, then $(d-1)x+dk\in\langle x\rangle+d\Delta=S$. Conversely, assume
\((d-1)x+dk\in S\). Then necessarily \((d-1)x+dk\in\mathbb{N}\), and $(d-1)x+dk=ax+d\delta$ for some $a\in\mathbb{N}$ and
$\delta\in\Delta$. Reducing modulo $d$ yields $ax\equiv -x\pmod d$, hence
$a\equiv d-1\pmod d$ since $\gcd(x,d)=1$. Writing $a=d-1+dt$ with
$t\in\mathbb{N}$ and dividing by $d$ gives $k=tx+\delta\in\Delta$.

Fix $f\in\operatorname{PF}(\Delta)$ and set $w=df+(d-1)x$. Since $f\notin\Delta$,
the equivalence implies $w\notin S$. Moreover, since $x\in\Delta\setminus\{0\}$,
the defining property of pseudo-Frobenius numbers gives $f+x\in\Delta$, and therefore
$w+x=d(f+x)\in d\Delta\subseteq S$.

Let $s\in S\setminus\{0\}$. If $s\in d\Delta$, say $s=d\delta$ with
$\delta\in\Delta\setminus\{0\}$, then $f+\delta\in\Delta$ and
$w+s=(d-1)x+d(f+\delta)\in S$. If $s\notin d\Delta$, write
$s=qx+d\delta$ with $q\ge 1$ and $\delta\in\Delta$. Then
$s-x=(q-1)x+d\delta\in S$, so $w+s=(w+x)+(s-x)\in S$. Hence
$w\in\operatorname{PF}(S)$, proving
\[
\{df+(d-1)x:\ f\in\operatorname{PF}(\Delta)\}\subseteq \operatorname{PF}(S).
\]

Conversely, let $w\in\operatorname{PF}(S)$. Since $x\in S\setminus\{0\}$, we have
$w+x\in S$, so $w+x=ax+d\delta$ for some $a\in\mathbb{N}$ and $\delta\in\Delta$.
If $a\ge 1$, then $w=(a-1)x+d\delta\in S$, a contradiction. Hence $a=0$ and
$w+x=d\delta$ for some $\delta\in\Delta$. Write $\delta=x+f$ with
$f\in\mathbb{Z}$, so $w=(d-1)x+df$.

If $f\in\Delta$, then $w\in S$ by the equivalence, again a contradiction; thus
$f\notin\Delta$. Finally, for every $\delta'\in\Delta\setminus\{0\}$ we have
$w+d\delta'\in S$, and
$w+d\delta'=(d-1)x+d(f+\delta')$. Using the equivalence, this forces
$f+\delta'\in\Delta$ for all $\delta'\in\Delta\setminus\{0\}$, hence
$f\in\operatorname{PF}(\Delta)$.

Therefore
\[
\operatorname{PF}(S)=\{df+(d-1)x:\ f\in \operatorname{PF}(\Delta)\}.
\]
The map $f\mapsto df+(d-1)x$ is a bijection, so $\operatorname{t}(S)=\operatorname{t}(\Delta)$.

Finally, a numerical semigroup is symmetric if and only if its type equals $1$,
equivalently $\operatorname{PF}(S)=\{\operatorname{F}(S)\}$. Hence $S$ is symmetric if and only if $\Delta$ is symmetric. 
\end{proof}

\end{document}
